\section{典型场景}

下面四个场景都来自漫展/晚会期间真实发生过的尴尬瞬间。每个场景说一下平时怎么走、用\,\app\,会怎么走。

\subsection{场景 A · 漫展开/闭幕式}

\textbf{现状。} 主持人在台前串场,PPT 跑在后台机器上由工作人员盯翻页。每次主持人说"下一页",工作人员要等 0.5\,$\sim$\,1.5\,s 才确认动作,节奏被拖住;主持人讲一段又得回头看后台是否跟上。

\textbf{用\,\app。} 主持人手机扫码,ARM 之后用一个全屏大按钮翻页(底栏占满整个手指能到的区域,单手大拇指可按)。下一页延迟 $\sim$\,30-80\,ms,主持人不再回头看后台。后台导播可同时看到当前页 + 状态,有问题随时\role{紧急锁}。

\textbf{典型流程。} 开幕前导播打开远控 → 主持人扫码 ARM → 开幕进行中只翻页 → 节目穿插期间不动 → 结束后导播\role{关远控}。

\subsection{场景 B · 互动答题 / 投票环节}

\textbf{现状。} PPT 内嵌"四选一"按钮(VBA 宏触发)。一般翻页器只能"上/下页",触发不了按钮,得有工作人员守着鼠标,按主持人口令操作。出错就是公开尴尬。

\textbf{用\,\app。} \app\ 自动检测当前幻灯片上所有绑定了\,\texttt{ppActionRunMacro}\,的形状,把它们的名字推到主持人页 —— 比如 "选项 A" / "选项 B" / "公布答案" 直接出现在主持人界面上,2 列大按钮排好。主持人在台前自己点哪个就触发哪个,工作人员只需要监看不需要操作。

\begin{infobox}
\textbf{技术细节}:形状 → 宏 的绑定关系由 \app\ 服务端实时从 PowerPoint 读出,主持人端只能看到\textit{形状名字}发回服务端,真正的宏调用由服务端再次查询 \texttt{ActionSettings(ppMouseClick).Run} 完成。即使主持人页被监听,也无法借此触发任意宏。
\end{infobox}

\subsection{场景 C · cosplay 大赛评分公布}

\textbf{现状。} 评委评分结束,主持人在台上准备公布前几名。后台工作人员要在 PPT 上一页页翻、找到对应排名页、按节奏推进。评分到颁奖的几分钟里,工作人员压力陡增,容易翻错页。

\textbf{用\,\app。} 评分 PPT 提前准备好按排名分页(第 1 名、第 2 名、第 3 名 …)。主持人在台上每念一个名字之前自己翻到对应页;后台导播全程只盯画面是否正确,出错可立即\role{回到上一页}或\role{黑屏}重整节奏。

\subsection{场景 D · 多段节目串场报幕}

\textbf{现状。} 漫展通常 10$\sim$15 个节目穿插,每个节目前要报幕、念赞助、放预告片。报幕员经常是不同的人,工作机和翻页器不能跟着报幕员换手。

\textbf{用\,\app。} 不同的报幕员各自扫一次码,导播这边可随时\role{收回 ARM}让控制权移交给下一位主持人。整场只有一台后台机器,不需要切设备;扫码一次的人不需要装任何 App,演出结束直接关浏览器走人。

\begin{infobox}
\textbf{操作模型}:任意时刻只有最近一次 ARM 的主持人能下指令,导播 GUI 里也能看到当前主持人的简短标识。"上一个报幕员还没正式交接、下一个就抢动作"这种情况不会发生 —— 不是靠人沟通靠住的,是流程上就不让它发生。
\end{infobox}
